% 0.摘要.tex —— 摘要
\begin{center}
  {\LARGE\bfseries 药材烘干过程温度与水分浓度分布的建模与数值求解}\\[6pt]
  {\large ——基于节点分类离散与线性化隐式差分的一维热质传递模型}
\end{center}

\begin{abstract}
\noindent 药材干燥是中药材品质控制的关键工序，预热平衡阶段内部温度与水分浓度的时空演化规律直接影响后续恒温干燥的工艺设定。本文针对给定半径 $2$ cm、长 $25$ cm 的圆柱形药材，建立径向一维非稳态热质传递模型，求解前 $30$ min 内药材内部温度与含水率的时空分布。

对于问题一，首先依据长径比与特征扩散长度将问题降维为无限长圆柱的径向一维问题，建立热传导方程与水分扩散方程；其次\textbf{完整实现 Euler 显式（向前 Euler）方法}，把径向节点严格划分为\textbf{中心点、内部点、边界点}三类分别离散：中心点用 L'H\^opital 法则消去 $1/r$ 奇异项并结合轴对称条件，内部点用柱坐标守恒型中心差分，边界点用 Robin 对流边界的半控制体积离散；显式格式逐点推进、无需解线性方程组、对变扩散系数 $D(C)$ 亦无需迭代，但受抛物型 CFL 稳定条件 $Fo=\alpha\Delta t/\Delta r^2\le0.41$ 限制（柱坐标下略严于笛卡尔情形的 $1/2$）。在此基础上\textbf{引入 Crank--Nicolson 隐式格式与 Thomas 追赶法}，将时间精度由一阶提升至二阶、并解除 CFL 限制实现无条件稳定，三对角方程组由 Thomas 追赶法以 $O(N)$ 复杂度求解。

结果表明，热扩散系数约为水分扩散系数的 $34$ 倍（$\mathrm{Fo}_T=0.760$、$\mathrm{Fo}_C=0.0222$），$1800$ s 时温度场已整体抬升（中心 $33.575$ $^\circ$C、表面 $36.785$ $^\circ$C、径向温差 $3.21$ $^\circ$C），而水分场仅外层约 $1.5$ cm 壳体失水（表面含水率降幅 $40.70\%$、体积平均降幅 $10.11\%$）。数值验证表明：Euler 显式与 CN 在 $\Delta t=1$ s 下四位小数一致；CN 格式可放大时间步 $300$ 倍而不失稳，显式格式在 $Fo>0.41$ 时迅速发散。

对于问题二，针对持续 $2\sim3$ 天的整个烘干过程（预热平衡 + 恒温干燥），建立\textbf{非线性变参数热质传递模型}：密度、比热容、导热系数随含水率变化，水分扩散系数同时随含水率与温度变化（附录 3）。数值上采用\textbf{线性化隐式 FDM + 迭代}：将 $D(C,T)$ 在当前状态处作一阶 Taylor 展开线性化（Newton 型，含 $\partial D/\partial C$ 敏度项），每个时间步内 Picard 迭代 $2\sim3$ 次即收敛，温度与水分方程以\textbf{弱耦合交替}方式求解、由 Thomas 追赶法完成三对角求解。结果表明：药材约 $2$ h 完成预热并进入恒温干燥，表面含水率数小时逼近环境平衡值，中心含水率降至 $0.15$ kg/kg 约需 $57$ h（约 $2.4$ 天），与题设“$2\sim3$ 天”吻合。此外给出问题三模型的推荐：沿用弱耦合交替求解并引入\textbf{自适应时间步}（按变化率调节 $\Delta t$），可将长时间求解成本压缩约两个数量级。

\noindent\textbf{关键词：}径向一维热质传递；非线性变参数模型；Taylor 线性化；Picard 迭代；Crank--Nicolson 格式；Thomas 追赶法；自适应时间步
\end{abstract}
